% sec-03: quantitative evidence, ARPA framing (what is already de-risked).
\section{What Is Already De-Risked}

An ARPA program should not spend its first year discovering whether the
hypothesis is worth testing. Three questions that would normally occupy that
year are answered and published.

\textbf{Is the drug worth the operation?} An August 2025 quantitative systems
pharmacology simulation, ten arms and 250 ordinary differential equations per
patient, returned median overall survival of \textbf{12.8 months} for the
daraxonrasib combination against \textbf{5.4 months} for chemotherapy
\cite{qspmetpancre}. In May 2026 RASolute 302 reported \textbf{13.2} against
\textbf{6.6 months} in the RAS G12 population \cite{rasolute302}. The proximity
is chronology, not validation: 1000 simulated against 241 enrolled, a
combination against a single agent, and KRAS G12C against a primarily G12D and
G12V population.

\textbf{Can the modeling be audited?} The digital-twin work carries an ASME
V\&V 40 and ICH M15 aligned credibility score of \textbf{81.9} across 55
verification notebook tests \cite{daraxsims,asmevv40,ichm15}.

\begin{apptable}
\begin{tabularx}{\textwidth}{>{\raggedright\arraybackslash}p{3.6cm} >{\raggedright\arraybackslash}p{2.6cm} >{\raggedright\arraybackslash}X}
\headrow \thead{De-risked item} & \thead{Result} & \thead{Why it matters at a gate}\\
\midrule
Arm selection & mOS HR 0.25, best arm & Gate 2 does not have to re-derive the regimen \\
Model credibility & Score 81.9, 55 tests & Gate 1 inherits an audited model rather than a new one \\
Regulatory package & IND drafted, protocol published & Gate 1 is an activation check, not an authoring task \\
Cost per simulation & \$36,330 against a \$120,000 to \$2,000,000 benchmark & The performer's burn rate is already demonstrated \\
\end{tabularx}
\end{apptable}
\tabcap{Four items retired before any federal dollar was requested, and the gate
each one removes work from. Every figure is drawn from the author's published
simulation and regulatory sources.}
